\section{Raw missing mass times native fiber degree}
\label{sec:tpc60-mass-degree}

We now give a phase-blind sufficient condition.  It is sharp at the level
of information it uses.

\subsection{The exact grouping norm}

Let \(\cW_M\) be a finite missing atomic carrier, let \(\cI_{\rm out}\) be
the active native output set, and let
\begin{equation}
 \cF_i^{\rm miss}
 :=\{w\in\cW_M:i(w)=i\}.
 \label{eq:tpc60-missing-native-fibers}
\end{equation}
Absorb every fixed unit M\"obius or high-character phase into the atomic
amplitude \(m_w(z)\), and define
\begin{align}
 L_M:\cK&\longrightarrow\ell^2(\cW_M),
 &(L_Mz)_w&:=m_w(z),
 \label{eq:tpc60-missing-atomic-lift}\\
 C_M:\ell^2(\cW_M)&\longrightarrow\ell^2(\cI_{\rm out}),
 &(C_Ma)_i&:=\sum_{w\in\cF_i^{\rm miss}}a_w.
 \label{eq:tpc60-missing-grouping}
\end{align}
Then
\begin{equation}
 M=C_ML_M,
 \qquad
 \cD_{\rm miss}(z):=\|L_Mz\|_2^2.
 \label{eq:tpc60-M-factorization-Dmiss}
\end{equation}
In the fixed-shift specialization, every absorbed phase has modulus one,
so \(\cD_{\rm miss}(z)=\cD_M(z)\) with the raw diagonal of
\eqref{eq:tpc60-BRM-diagonals}.
Put
\begin{equation}
 d_{\rm miss}:=\max_{i\in\cI_{\rm out}}
                   |\cF_i^{\rm miss}|,
 \label{eq:tpc60-native-missing-degree}
\end{equation}
with \(d_{\rm miss}=0\) when the missing carrier is empty.

\begin{proposition}[Exact native grouping norm]
\label{prop:tpc60-exact-grouping-norm}
One has
\begin{equation}
 \boxed{
 \|C_M\|^2=d_{\rm miss},}
 \label{eq:tpc60-CM-norm}
\end{equation}
and therefore, for every \(z\in\cK\),
\begin{equation}
 \boxed{
 \|Mz\|^2\le d_{\rm miss}\cD_{\rm miss}(z).}
 \label{eq:tpc60-mass-degree-upper}
\end{equation}
\end{proposition}

\begin{proof}
Fiberwise Cauchy--Schwarz gives
\[
 |(C_Ma)_i|^2
 \le|\cF_i^{\rm miss}|
    \sum_{w\in\cF_i^{\rm miss}}|a_w|^2.
\]
Summing over the disjoint output fibers proves
\(\|C_M\|^2\le d_{\rm miss}\).  On a fiber of maximal size, choose all
coordinates equal with a common phase and set all other coordinates to
zero.  Equality follows.  Applying the operator norm bound to \(L_Mz\)
proves \eqref{eq:tpc60-mass-degree-upper}.
\end{proof}

The exact norm concerns the ambient grouping map.  The physical range
\(\Ran L_M\) may have a smaller restricted norm, but proving that is another
Gram problem.  Conversely, replacing \(d_{\rm miss}\) by an average degree
is invalid for a coefficient-uniform estimate.

\subsection{Distinguished-vector certificate}

\begin{theorem}[Sharp mass--degree lower certificate]
\label{thm:tpc60-pointwise-mass-degree}
For every \(z\) with \(Az\ne0\),
\begin{equation}
 \boxed{
 \frac{\|(A+M)z\|^2}{\|Az\|^2}
 \ge
 \left(
  1-
  \sqrt{\frac{d_{\rm miss}\cD_{\rm miss}(z)}{\|Az\|^2}}
 \right)_+^2.}
 \label{eq:tpc60-pointwise-mass-degree-bound}
\end{equation}
In particular, if
\begin{equation}
 d_{\rm miss}\cD_{\rm miss}(z)
 \le\eta\|Az\|^2,
 \qquad 0\le\eta<1,
 \label{eq:tpc60-eta-mass-degree}
\end{equation}
then the ratio is at least \((1-\sqrt\eta)^2\).
\end{theorem}

\begin{proof}
The reverse triangle inequality and
\cref{prop:tpc60-exact-grouping-norm} give
\[
 \|(A+M)z\|
 \ge(\|Az\|-\|Mz\|)_+
 \ge
 \left(
  \|Az\|-\sqrt{d_{\rm miss}\cD_{\rm miss}(z)}
 \right)_+.
\]
Divide by \(\|Az\|\) and square.
\end{proof}

\begin{example}[Sharp cancellation at the product threshold]
\label{ex:tpc60-sharp-mass-degree-threshold}
Take one represented native output \(Az=1\).  Let one native fiber contain
\(n\) missing atoms, each of amplitude \(-1/n\).  Then
\begin{equation}
 d_{\rm miss}=n,
 \qquad
 \cD_{\rm miss}(z)=\frac1n,
 \qquad
 Mz=-1.
 \label{eq:tpc60-threshold-example-data}
\end{equation}
Thus \(d_{\rm miss}\cD_{\rm miss}/\|Az\|^2=1\) and the full output
vanishes.  The same construction with \(n=2\) shows that very small
combinatorial degree is already compatible with exact cancellation.
Therefore neither \(\cD_{\rm miss}=o(1)\) nor a bounded degree alone is a
sufficient reassembly theorem.
\end{example}

\subsection{Uniform and canonical-sector versions}

\begin{corollary}[Uniform mass--degree certificate]
\label{cor:tpc60-uniform-mass-degree}
Suppose \(\varepsilon\ge0\) and
\begin{equation}
 \cD_{\rm miss}(z)\le\varepsilon\|Az\|^2
 \qquad\text{for every }z\in\cK.
 \label{eq:tpc60-uniform-missing-mass-hypothesis}
\end{equation}
Then \(M|_{\Ker A}=0\) and
\begin{equation}
 \boxed{
 \delta(A,M)
 \ge(1-\sqrt{d_{\rm miss}\varepsilon})_+^2.}
 \label{eq:tpc60-uniform-mass-degree-bound}
\end{equation}
\end{corollary}

\begin{proof}
If \(z\in\Ker A\), \eqref{eq:tpc60-uniform-missing-mass-hypothesis}
forces \(\|L_Mz\|=0\), hence \(Mz=0\).  Apply
\cref{thm:tpc60-pointwise-mass-degree} to every \(z\) with \(Az\ne0\) and
take the infimum.
\end{proof}

The hypothesis in \cref{cor:tpc60-uniform-mass-degree} is strong: it
automatically removes the nuisance image.  When only the canonical sector
can be controlled, the following version retains the missing
transversality explicitly.  With the notation of
\cref{sec:tpc60-shorted-gram}, define
\begin{equation}
 \gamma_A
 :=\inf_{s\in S\setminus\{0\}}
   \frac{\|PA_Ss\|^2}{\|A_Ss\|^2}.
 \label{eq:tpc60-gamma-A}
\end{equation}

\begin{proposition}[Canonical mass plus nuisance transversality]
\label{prop:tpc60-canonical-mass-transversality}
Suppose \(\varepsilon\ge0\) and
\begin{equation}
 \cD_{\rm miss}(s)
 \le\varepsilon\|A_Ss\|^2
 \qquad(s\in S).
 \label{eq:tpc60-canonical-missing-bound}
\end{equation}
Then
\begin{equation}
 \boxed{
 \delta(A,M)
 \ge
 \left(
  \sqrt{\gamma_A}-\sqrt{d_{\rm miss}\varepsilon}
 \right)_+^2.}
 \label{eq:tpc60-canonical-transversality-bound}
\end{equation}
\end{proposition}

\begin{proof}
For \(s\in S\),
\begin{align*}
 \|PB_Ss\|
 &=\|P(A_Ss+Ms)\|\\
 &\ge\|PA_Ss\|-\|PMs\|\\
 &\ge
 \left(\sqrt{\gamma_A}-\sqrt{d_{\rm miss}\varepsilon}\right)
 \|A_Ss\|.
\end{align*}
Use \cref{thm:tpc60-exact-nuisance-elimination} and take the positive part.
\end{proof}

This estimate exposes two independent risks: the missing atomic output can
be large, and the represented output can already lie close to the nuisance
range \(B(\Ker A)\).  An estimate for only one of them does not close the
uniform gate.

\subsection{Exponent form}

For a distinguished vector \(z_*(X)\), write
\begin{equation}
 \frac{\|B_Xz_*(X)\|^2}{\|A_Xz_*(X)\|^2}
 \ge X^{-\lambda_{\rm reasm,*}+o(1)}
 \label{eq:tpc60-pointwise-reassembly-exponent}
\end{equation}
when the denominator is nonzero.  The star distinguishes this pointwise
exponent from the coefficient-uniform \(\lambda_{\rm reasm}\) of
\cref{sec:tpc60-kernel-tests}.  Suppose now, for that distinguished vector
or uniformly on a stated class,
\begin{equation}
 d_{\rm miss}=X^{\mu_{\rm miss}+o(1)},
 \qquad
 \frac{\cD_{\rm miss}(z)}{\|Az\|^2}
 \le X^{-\sigma_{\rm miss}+o(1)}.
 \label{eq:tpc60-degree-saving-exponents}
\end{equation}
Writing
\begin{equation}
 q_X(z):=
 \frac{d_{\rm miss}\cD_{\rm miss}(z)}{\|Az\|^2},
 \label{eq:tpc60-qX-product}
\end{equation}
the triangle inequality gives the two-sided estimate
\begin{equation}
 \bigl(1-\sqrt{q_X(z)}\bigr)_+^2
 \le\frac{\|Bz\|^2}{\|Az\|^2}
 \le\bigl(1+\sqrt{q_X(z)}\bigr)^2.
 \label{eq:tpc60-qX-two-sided}
\end{equation}
Consequently
\begin{equation}
 \sigma_{\rm miss}>\mu_{\rm miss}
 \quad\Longrightarrow\quad
 \frac{\|Bz\|^2}{\|Az\|^2}=1+o(1)
 \label{eq:tpc60-zero-cost-exponent-condition}
\end{equation}
Thus the mass--degree route gives \(\lambda_{\rm reasm,*}=0\) in the
distinguished-vector scope and \(\lambda_{\rm reasm}=0\) when the hypotheses
hold uniformly on the declared class.  At equality of the exponents, the exact
sufficient condition furnished by this certificate is
\begin{equation}
 \bigl(1-\sqrt{q_X(z)}\bigr)_+^2\ge X^{-o(1)}.
 \label{eq:tpc60-equality-subpower-test}
\end{equation}
This bound must hold pointwise for the distinguished vector, or uniformly over the
declared coefficient class, according to the claimed scope.
The simpler condition \(\limsup q_X(z)<1\) is sufficient but not
necessary; for example, \(q_X=1-1/\log X\) still gives a subpower lower
bound.  Thus exponent equality requires finer constant or subpower
control.  If
\(\sigma_{\rm miss}<\mu_{\rm miss}\), the certificate is silent; it does
not prove failure.

More generally, if a common reference diagonal satisfies
\begin{equation}
 \cD_{\rm miss}(z)
 \le X^{-\sigma+o(1)}\cD_{\rm ref}(z),
 \qquad
 \|Az\|^2
 \ge X^{-\lambda_{\rm rep}+o(1)}\cD_{\rm ref}(z),
 \label{eq:tpc60-common-reference-bounds}
\end{equation}
then the zero-cost sufficient condition becomes
\begin{equation}
 \boxed{
 \sigma>\mu_{\rm miss}+\lambda_{\rm rep}.}
 \label{eq:tpc60-common-reference-threshold}
\end{equation}
Every exponent in this comparison refers to the same physical vector and
the same raw reference measure.
